% =====================================================================
% P9 — Threshold collapse of the inverted-U under extreme institutional
%       turbulence: India 2014-2025 (WBES three waves)
% WORKING DRAFT submission framework (Springer-compatible).
%
% COMPILE:
%   pdflatex p9_india_threshold.tex (twice) ; bibtex p9_india_threshold ; pdflatex (x2)
% NOTE: No pdflatex was available in the build environment. Uses standard
%   `article` class (Springer sn-jnl.cls not bundled). Swap class/title for
%   sn-jnl when submitting to a Springer journal (e.g. MIR).
%
% INTEGRITY NOTE: Every numeric value is verbatim from the dissertation
%   (chuong_4_ket_qua_vi.md sec 4.8 / 4.9, incl. Table 4.3; methods from
%   phu_luc_A). Full per-wave coefficient tables absent from the dissertation
%   are explicit PLACEHOLDERS. No new coefficients invented.
% =====================================================================

\documentclass[12pt,a4paper]{article}
\usepackage[margin=2.5cm]{geometry}
\usepackage{setspace}\onehalfspacing
\usepackage{iftex}
\ifPDFTeX
  \usepackage[utf8]{inputenc}\usepackage[T1]{fontenc}\usepackage{lmodern}
\else
  \usepackage{fontspec}
\fi
\usepackage{amsmath,amssymb,amsfonts,mathtools}
\usepackage{booktabs}\usepackage{array}\usepackage{tabularx}
\usepackage{graphicx}\usepackage{float}
\usepackage{caption}\captionsetup{font=small,labelfont=bf}
\usepackage{microtype}\usepackage{textcomp}
\usepackage[hidelinks,breaklinks=true]{hyperref}
\usepackage[round]{natbib}\bibliographystyle{apalike}

\usepackage{filecontents}
\begin{filecontents*}{p9refs.bib}
@article{deaton1985,
  author={Deaton, Angus}, title={Panel data from a time series of cross-sections},
  journal={Journal of Econometrics}, volume={30}, number={1--2}, pages={109--126}, year={1985}}
@incollection{verbeek2008,
  author={Verbeek, Marno}, title={Pseudo-panels and repeated cross-sections},
  booktitle={The Econometrics of Panel Data}, edition={3rd}, pages={369--383},
  publisher={Springer}, year={2008}}
@article{lindmehlum2010,
  author={Lind, Jo Thori and Mehlum, Halvor},
  title={With or without {U}? The appropriate test for a {U}-shaped relationship},
  journal={Oxford Bulletin of Economics and Statistics}, volume={72}, number={1},
  pages={109--118}, year={2010}}
@article{paternoster1998,
  author={Paternoster, Raymond and Brame, Robert and Mazerolle, Paul and Piquero, Alex},
  title={Using the correct statistical test for the equality of regression coefficients},
  journal={Criminology}, volume={36}, number={4}, pages={859--866}, year={1998}}
@article{mizik2003,
  author={Mizik, Natalie and Jacobson, Robert},
  title={Trading off between value creation and value appropriation: The financial implications
  of shifts in strategic emphasis},
  journal={Journal of Marketing}, volume={67}, number={1}, pages={63--76}, year={2003}}
@article{hsieh2009,
  author={Hsieh, Chang-Tai and Klenow, Peter J.},
  title={Misallocation and manufacturing {TFP} in {China} and {India}},
  journal={Quarterly Journal of Economics}, volume={124}, number={4}, pages={1403--1448}, year={2009}}
@book{north1990,
  author={North, Douglass C.}, title={Institutions, Institutional Change and Economic Performance},
  publisher={Cambridge University Press}, year={1990}}
@book{khanna2010,
  author={Khanna, Tarun and Palepu, Krishna G.},
  title={Winning in Emerging Markets: A Road Map for Strategy and Execution},
  publisher={Harvard Business Press}, year={2010}}
@article{cameron2015,
  author={Cameron, A. Colin and Miller, Douglas L.},
  title={A practitioner's guide to cluster-robust inference},
  journal={Journal of Human Resources}, volume={50}, number={2}, pages={317--372}, year={2015}}
@incollection{dophan2025,
  author={Do, T. H. and Phan, A. T.},
  title={Internationalization and firm performance of firms in {India}: The role of top management},
  booktitle={International Business Research: Traditional and Creative Approaches},
  editor={Bartekova, M.}, publisher={IntechOpen}, year={2025},
  note={https://doi.org/10.5772/intechopen.1011012}}
\end{filecontents*}

\title{Threshold Collapse: How Extreme Institutional Turbulence Dissolves the
Internationalization--Performance Inverted-U in India, 2014--2025}
\author{Author details withheld for blind review}
\date{Working draft}

\begin{document}
\maketitle

\begin{abstract}
\noindent Whether the inverted-U turning point of the
internationalization--performance (I--P) relationship is structurally stable
under institutional change is unknown. We exploit India's exceptional
2014--2025 reform sequence (demonetization, the Insolvency and Bankruptcy
Code, the Goods and Services Tax, the Unified Payments Interface, and
Make-in-India/Atmanirbhar Bharat) as a quasi-experiment, using three World
Bank Enterprise Survey waves (2014 PICS3, $N=8{,}941$; 2022 BEE, $N=9{,}300$;
2025 BREADY, $N=10{,}476$; pooled $N=28{,}717$). In 2014 the relationship is a
classic inverted-U ($\beta_1=+1.865$, $\beta_2=-1.508$, both $p<0.001$;
turning point $61.8\%$ of export intensity). By 2022 the curve persists but
its turning point contracts inward to $40.7\%$. By 2025 the quadratic term is
insignificant ($\beta_2=-0.160$, $p=0.420$) and the relationship is
monotonically negative---the threshold has dissolved. A Paternoster test
rejects coefficient equality between 2014 and 2025 with extreme strength
($z_{\mathrm{FSTS}}=-7.94$, $z_{\mathrm{FSTS}^2}=+4.17$; both $p<0.0001$;
state-clustered $z=-3.50$/$+2.19$). The pattern contrasts sharply with China,
where the turning point is structurally stable across twelve years. We
interpret threshold collapse as evidence that I--P turning-point stability is
not an intrinsic property of an institutional regime but depends on the pace
and depth of cumulative institutional change.\\[4pt]
\textbf{Keywords:} internationalization; firm performance; inverted-U;
threshold stability; institutional turbulence; India; digital public
infrastructure
\end{abstract}

\section{Introduction}
The location of the I--P turning point is typically treated as a fixed
structural parameter of a context. Whether that parameter survives rapid
institutional change has not been tested, because few settings combine a
large firm-level panel of cross-sections with a concentrated burst of reform.
India 2014--2025 provides exactly that. The gap we address is the temporal
\emph{stability} of the inverted-U threshold under extreme institutional
turbulence. Our contribution is to document \emph{threshold collapse}: a
turning point that contracts and then dissolves as cumulative institutional
shocks accumulate, and to contrast it with the structurally stable Chinese
case, showing that stability is conditional rather than regime-intrinsic. The
study extends prior published work on India \citep{dophan2025} from a single
cross-section to a three-wave stability test.

\section{Theory and Hypotheses}
Institutional theory holds that the costs and benefits of internationalization
are conditioned by the institutional environment \citep{north1990,khanna2010}.
If that environment is itself reconfigured---through monetary, tax, insolvency,
and digital-infrastructure reform---the cost--benefit balance underlying the
inverted-U should shift, relocating and potentially erasing the turning point.
Where reforms raise the relative payoff of domestic orientation (for example
production-linked incentives and self-reliance policy), the ascending arm of
the export curve should weaken over time.

\begin{description}
\item[H1c.] Under extreme institutional turbulence the inverted-U turning
point shifts and ultimately dissolves over time (threshold collapse).
\end{description}

\section{Data and Methods}
The data are three WBES India waves treated as repeated cross-sections
\citep{deaton1985,verbeek2008}: 2014 PICS3 ($N=8{,}941$), 2022 BEE
($N=9{,}300$), and 2025 BREADY ($N=10{,}476$), pooled $N=28{,}717$. [Sample
sizes are REAL.] The dependent variable is within-wave standardized labour
productivity; FSTS is export intensity. Each wave is fitted with the quadratic
specification and assessed with the Lind--Mehlum test \citep{lindmehlum2010};
turning points are $-\beta_1/(2\beta_2)$ with delta-method confidence
intervals. Cross-wave coefficient equality is tested with the Paternoster
$z$ test \citep{paternoster1998} under both HC1 and state-clustered
($G=24$ states) standard errors \citep{cameron2015}. A pooled
specification with wave dummies and interactions confirms the systematic
shift.

\section{Results}
\subsection{From classic inverted-U to collapse (H1c)}
In 2014 the relationship is a classic inverted-U:
$\beta_1=+1.865$ ($p<0.001$), $\beta_2=-1.508$ ($p<0.001$), turning point
$\mathrm{TP}=61.8\%$ FSTS (delta-method $95\%$ CI $[55.0\%, 68.6\%]$),
Lind--Mehlum $p<0.001$. By 2022, after the first reform cluster, the curve
persists but its turning point contracts inward: $\beta_1=+1.542$,
$\beta_2=-1.893$, $\mathrm{TP}=40.7\%$ (CI $[38.0\%, 43.5\%]$). By 2025 the
quadratic term loses significance ($\beta_2=-0.160$, $p=0.420$) and the
relationship is monotonically negative over the observed export range. [All
values in this subsection are REAL.]

\begin{table}[H]
\centering
\caption{Turning point by survey wave, India (REAL values from dissertation
Table~4.3).}
\begin{tabularx}{\textwidth}{lccX}
\toprule
Wave & $N$ (analytic) & Turning point & Lind--Mehlum \\
\midrule
2014 (PICS3)  & $8{,}941$  & $61.8\%$ FSTS & $p<0.001$ (classic inverted-U) \\
2022 (BEE)    & $9{,}300$  & $40.7\%$ FSTS & $p<0.001$ (TP shifts inward) \\
2025 (BREADY) & $10{,}476$ & undefined     & $p=0.99$ (curve collapses) \\
\bottomrule
\end{tabularx}
\end{table}

\subsection{Cross-wave stability tests}
The Paternoster test rejects 2014--2025 coefficient equality with extreme
strength: $z(\mathrm{FSTS})=-7.94$, $p<0.0001$ and
$z(\mathrm{FSTS}^2)=+4.17$, $p<0.0001$ under HC1; under state-clustered
standard errors ($G=24$) the result remains significant,
$z(\mathrm{FSTS})=-3.50$ ($p=0.001$) and $z(\mathrm{FSTS}^2)=+2.19$
($p=0.029$). The pooled specification with wave interactions confirms the
shift, $\beta(\mathrm{FSTS}\times\mathrm{wave}\_2025)=-1.63$ ($p<0.0001$).
The contrast with China is direct: there the turning point is structurally
stable ($\mathrm{TP}=49.4\%$ in 2012 vs.\ $47.2\%$ in 2024; Paternoster
$z=+0.82$, $p=0.412$). [All values in this subsection are REAL.]

% --------------------------------------------------------------------
\begin{table}[H]
\centering
\caption{Per-wave quadratic regressions with TCI and DAI moderators
(full coefficients).}
\begin{tabularx}{\textwidth}{Xccc}
\toprule
 & 2014 & 2022 & 2025 \\
\midrule
\multicolumn{4}{c}{\itshape [Table P9-2 -- full coefficients from the authors'
original output. Populate only from locked estimates.]}\\
FSTS         & $+1.865^{***}$ & $+1.542$ & --- \\
FSTS$^2$     & $-1.508^{***}$ & $-1.893$ & $-0.160$ (ns) \\
TCI (direct) & $+0.12^{***}$  & $+0.19^{***}$ & $-0.08^{***}$ \\
FSTS$\times$TCI & --- & --- & $+0.20\ (p=0.03)$ \\
FSTS$\times$DAI$_{\text{Tier-2}}$ & --- & --- & $-4.02\ (p=0.004)$ \\
$N$          & $8{,}941$ & $9{,}300$ & $10{,}476$ \\
\bottomrule
\end{tabularx}
\end{table}
% --------------------------------------------------------------------

\subsection{Capability and digital moderators}
TCI shows a sign flip unobserved elsewhere: positive in 2014
($\beta=+0.12$, $p<0.001$) and stronger in 2022 ($\beta=+0.19$, $p<0.001$),
but negative in 2025 ($\beta=-0.08$, $p<0.001$) with a positive
FSTS$\times$TCI interaction ($\beta=+0.20$, $p=0.03$), consistent with a
strategic-emphasis trade-off \citep{mizik2003} as post-2020 industrial policy
redirected capability toward domestic production. For digital adoption, the
2025 BREADY wave adds a Tier-2 indicator (electronic-payment revenue share),
whose interaction with export intensity is negative and significant
($\mathrm{FSTS}\times\mathrm{DAI}_{\text{Tier-2}}=-4.02$, $p=0.004$):
domestic-rail public digital infrastructure (UPI) does not substitute for the
cross-border financial infrastructure exporters require. [All values in this
subsection are REAL.]

\section{Discussion and Conclusion}
The Indian evidence shows that the I--P turning point is not an immutable
structural constant. Under a concentrated sequence of monetary, fiscal,
insolvency, and digital reforms, the inverted-U first contracted and then
dissolved into a monotone-negative relationship. The contrast with the
structurally stable Chinese case---both upper-middle-income regimes---implies
that threshold stability is governed by the pace and depth of cumulative
institutional change within the observation window, not by the institutional
regime label alone. The TCI sign flip and the negative Tier-2 digital
interaction further indicate that capability and digital effects are
themselves policy-contingent: public digital infrastructure substitutes for
private capability only when the two share a purpose orientation.

Limitations include the repeated-cross-section design, which precludes firm
trajectory tracking, and the inability to attribute collapse to any single
reform within a bundled sequence. The findings nonetheless caution against
treating any estimated turning point as a fixed planning parameter in rapidly
reforming economies. An extended manuscript testing I--P threshold robustness
across the three 2014--2025 waves is in preparation for MIR/IJOEM.

\bibliography{p9refs}

\end{document}
